% !TeX root = ../main.tex
\section{задача 3-1}

\vspace{1em}
\textbf{1. Анализ исходных данных}

Согласно условию, отношение сигнал/шум на входе детектора велико ($a \gg 3$). Это означает, что частотный детектор (ЧД) работает в линейном (надпороговом) режиме, и подавления сигнала шумом не происходит.

Как показано на структурной схеме, спектральная плотность мощности (СПМ) шума на выходе ЧД имеет параболическую зависимость от частоты и определяется формулой:
\begin{equation*}
G_{\text{ЧД}}(F) \approx S_{\text{ЧД}}^2 \frac{F^2}{U_c^2} \cdot 2K_{0 \text{ БВЧ}}^2 G_0
\end{equation*}
Для удобства вычислений обозначим независящую от частоты $F$ часть выражения как константу $A$:
\begin{equation*}
A = S_{\text{ЧД}}^2 \frac{2K_{0 \text{ БВЧ}}^2 G_0}{U_c^2} \quad \Rightarrow \quad G_{\text{ЧД}}(F) = A \cdot F^2
\end{equation*}

Мощность (дисперсия) шума на выходе идеального прямоугольного фильтра БНЧ с полосой пропускания $F_{\text{в}}$ рассчитывается как интеграл от СПМ:
\begin{equation*}
U_{\text{ш.вых}}^2 = \int_0^{F_{\text{в}}} G_{\text{ЧД}}(F) dF = \int_0^{F_{\text{в}}} A \cdot F^2 dF = A \left[ \frac{F^3}{3} \right]_0^{F_{\text{в}}} = \frac{A}{3} F_{\text{в}}^3
\end{equation*}

Среднеквадратическое (эффективное) напряжение шума пропорционально квадратному корню из его мощности:
\begin{equation*}
U_{\text{эфф}} = \sqrt{U_{\text{ш.вых}}^2} = \sqrt{\frac{A}{3}} \cdot F_{\text{в}}^{3/2}
\end{equation*}

\vspace{1em}
\textbf{2. Ответы на вопросы задачи}

\textbf{а) Изменение полосы пропускания БВЧ:}

Если полосу пропускания БВЧ увеличить вдвое, сохранив коэффициент усиления $K_{0 \text{ БВЧ}}$ неизменным, константа $A$ не поменяется. В надпороговом режиме ($a \gg 3$) спектральная плотность шума $G_{\text{ЧД}}(F)$ на выходе детектора в звуковом диапазоне \textbf{не зависит} от ширины полосы БВЧ. Высокочастотные составляющие шума, которые пропустит расширенный БВЧ, будут в дальнейшем отфильтрованы узкополосным фильтром БНЧ.

\textit{Вывод:} Эффективное напряжение шума на выходе БНЧ \textbf{не изменится}.

\vspace{1em}
\textbf{б) Изменение полосы пропускания БНЧ:}

Если полосу пропускания БНЧ расширить вдвое, то новая верхняя частота среза станет равна $F_{\text{в}}' = 2F_{\text{в}}$. 
Рассчитаем новую мощность шума на выходе:
\begin{equation*}
(U_{\text{ш.вых}}')^2 = \frac{A}{3} (2F_{\text{в}})^3 = \frac{A}{3} \cdot 8F_{\text{в}}^3 = 8 \cdot \left( \frac{A}{3} F_{\text{в}}^3 \right) = 8 \cdot U_{\text{ш.вых}}^2
\end{equation*}
Таким образом, мощность шума возрастет в 8 раз. 

Теперь найдем новое эффективное напряжение шума:
\begin{equation*}
U_{\text{эфф}}' = \sqrt{(U_{\text{ш.вых}}')^2} = \sqrt{8 \cdot U_{\text{ш.вых}}^2} = \sqrt{8} \cdot U_{\text{эфф}} = 2\sqrt{2} \cdot U_{\text{эфф}} \approx 2.83 \cdot U_{\text{эфф}}
\end{equation*}

\textit{Вывод:} Эффективное напряжение шума на выходе БНЧ \textbf{увеличится в $2\sqrt{2}$ (примерно в 2,83) раза}.